\documentclass[11pt,a4paper]{article}
%%%%%%%%%%%%%%%%%%%%%%%%%%%%%%%%%%%%%%%%%%%%%%%%%%%%%%%%%%%%%%%%%%%%%%%%%%%%%%
% Konzept Förderkurs Mathematik – Kurzfassung zur Diskussion
%%%%%%%%%%%%%%%%%%%%%%%%%%%%%%%%%%%%%%%%%%%%%%%%%%%%%%%%%%%%%%%%%%%%%%%%%%%%%%
\usepackage[T1]{fontenc}
\usepackage{lmodern}
\usepackage[nswissgerman]{babel}
\usepackage[margin=2.3cm]{geometry}
\usepackage{amsmath, amssymb}
\usepackage{booktabs, tabularx, array}
\usepackage[hidelinks]{hyperref}
\usepackage[most]{tcolorbox}
\usepackage{titlesec}
\usepackage{enumitem}
\usepackage{needspace}

\definecolor{fkblau}{RGB}{0,84,147}

\titleformat{\section}{\Large\bfseries\color{fkblau}}{\thesection}{0.6em}{}
\setlist{itemsep=0.3ex, topsep=0.5ex}
\setlength{\parindent}{0pt}
\setlength{\parskip}{0.8ex}
\renewcommand{\arraystretch}{1.35}
\newcolumntype{C}{>{\centering\arraybackslash}X}

\newcommand{\kurs}{Förderkurs Mathematik}

\begin{document}

\begin{center}
  {\LARGE\bfseries\color{fkblau} \kurs}\\[1ex]
  {\large Ideen für den Förderunterricht im ersten Jahr}\\[1ex]
  {\small Entwurf zur Diskussion · Stand Oktober 2026}
\end{center}

\begin{tcolorbox}[colframe=fkblau, colback=fkblau!5, boxrule=0.6pt, arc=2pt]
Dieses Papier fasst meine eigenen Überlegungen zum Förderunterricht
Mathematik zusammen. Es ist kein beschlossenes Konzept, sondern eine
Diskussionsgrundlage. Die beigelegten Beispiele zeigen, wie die Ideen konkret
aussehen könnten. Ich freue mich über Rückmeldungen, welche Teile überzeugen
und welche nicht.
\end{tcolorbox}

Viele Schwierigkeiten im Gymnasium liegen tiefer als der aktuelle Stoff: bei
Brüchen, beim Umgang mit Variablen oder beim Übersetzen von Text in
Mathematik. Der Kurs soll diese Lücken \textbf{gezielt} schliessen, statt
für alle dasselbe zu wiederholen.

% ===========================================================================
\section{Standortbestimmung zu Beginn}

Zu Beginn des Kurses lösen alle Teilnehmenden eine Standortbestimmung. Sie
zeigt nicht nur, \emph{dass} jemand Unterstützung braucht, sondern auch,
\emph{wobei}. Daraus ergibt sich für jede Schülerin und jeden Schüler ein
persönliches Profil, nach dem sie oder er im Kurs übt.

Ob die Standortbestimmung auch über die Teilnahme am Kurs entscheidet oder
ob dafür der Check S3 verwendet wird, ist noch offen.

% ===========================================================================
\section{Ein Raster in zwei Dimensionen}

Die Standortbestimmung und alle Übungen sind nach demselben Raster geordnet:
\textbf{Themen} (Stränge) in der einen Richtung, \textbf{Abstraktionsniveaus}
in der anderen. Am Beispiel der Brüche:

\begin{tabularx}{\textwidth}{@{}l*{4}{C}@{}}
  \toprule
  & \textbf{A Grundlagen} & \textbf{B Numerisch} & \textbf{C Symbolisch} & \textbf{D Formal} \\
  & \small Bedeutung, Zahl\-verständnis & \small Rechnen mit Zahlen
  & \small dasselbe mit Variablen & \small Parameter, Strukturen \\
  \midrule
  Brüche &
    $\frac{3}{4}$ auf dem Zahlenstrahl einzeichnen &
    $\dfrac{2}{3} + \dfrac{1}{4}$ &
    $\dfrac{a}{3} + \dfrac{b}{4}$ &
    $\dfrac{1}{x} + \dfrac{1}{x + 1}$ \\
  \bottomrule
\end{tabularx}

Die Stränge sind: Zahlen und Stellenwert, Operationen, Brüche,
Proportionalität und Prozent, Terme, Gleichungen, Funktionen und Graphen,
Geometrie und Grössen sowie Sprache und Mathematik.

Für jeden Strang zeigt die Standortbestimmung die \textbf{Bruchstelle}: das
Niveau, ab dem es nicht mehr sicher klappt. Wer $\frac{2}{3} + \frac{1}{4}$
richtig rechnet, bei $\frac{a}{3} + \frac{b}{4}$ aber scheitert, hat kein
Problem mit Brüchen, sondern mit dem Schritt zur Algebra. Und wer bei
Gleichungen mit Brüchen scheitert, hat oft gar kein Problem mit Gleichungen,
sondern eine Lücke bei den Brüchen.

% ===========================================================================
\needspace{12\baselineskip}
\section{Fortschritt sichtbar machen}

Im Verlauf des Jahres folgen weitere Standortbestimmungen in parallelen
Versionen (gleicher Aufbau, andere Zahlen), zum Beispiel:

\begin{center}
\begin{tabular}{llll}
  \textbf{August} & \textbf{Oktober/November} & \textbf{Januar} & \textbf{März} \\
  Ausgangsprofil & 1.\ Kontrolle & 2.\ Kontrolle & Abschluss
\end{tabular}
\end{center}

So sehen die Schülerinnen und Schüler, wie sich ihre Bruchstellen
verschieben, und die Lehrperson kann die Schwerpunkte laufend anpassen.
Sichtbarer Fortschritt ist zugleich der wirksamste Weg, Selbstvertrauen
aufzubauen.

% ===========================================================================
\section{Drei Pfade im Kurs}

\begin{tabularx}{\textwidth}{@{}>{\raggedright\arraybackslash}p{2.9cm}c X@{}}
  \toprule
  \textbf{Pfad} & \textbf{Anteil} & \textbf{Inhalt} \\
  \midrule
  Instruktion & ca.\ 20\,\% &
    Kurze \textbf{Mini-Lektionen} (ca.\ 10 Minuten) zu einer einzelnen
    Idee, oft nur für die Teilgruppe, die sie braucht: fehlende Konzepte,
    typische Fehlvorstellungen, saubere Darstellung. \\
  Üben & ca.\ 60\,\% &
    \textbf{Persönliche Übungsblätter} nach dem eigenen Profil, aus einer
    gemeinsamen Aufgabensammlung, die nach demselben Raster geordnet ist.
    Mit Lösungen zur Selbstkontrolle und Hinweisen zu typischen Fehlern. \\
  Werkzeugkasten & ca.\ 20\,\% &
    \textbf{Überfachliche Werkzeuge}: Texte in Mathematik übersetzen und
    umgekehrt, wirksam lernen, mit Fehlern und Prüfungsangst umgehen,
    Vertrauen in die eigenen Fähigkeiten. Kurz eingeführt und danach im
    Kurs immer wieder eingefordert. \\
  \bottomrule
\end{tabularx}

Der Kurs ist unabhängig vom Wochenstoff der einzelnen Klassen. Er behandelt
die Grundlagen in der Reihenfolge, in der sie im Regelunterricht gebraucht
werden: zuerst Algebra, dann Funktionen, dann quadratische Gleichungen und
Geometrie.

% ===========================================================================
\section*{Beilagen}

\begin{enumerate}
  \item Standortbestimmung mit Lösungen und Auswertungsraster
  \item Zwei Mini-Lektionen: «Das Gleichheitszeichen» und
        «Rechenreihenfolge und Klammern»
  \item Zwei Einheiten aus dem Werkzeugkasten: «Sprache und Mathematik» und
        «Vertrauen in die eigene Mathematik»
  \item Beispiele für Übungsblätter aus der Aufgabensammlung
\end{enumerate}

Alle Beilagen gibt es in einer Schüler- und einer Lehrpersonenversion.

\end{document}
